\section{Preliminaries}
\label{sec:prelim}

Throughout, all variables range over $\N=\{0,1,2,\dots\}$ unless stated
otherwise, and $\sq$ denotes ``is a perfect square''.

\begin{definition}[Representation]
A set $S\subseteq\N$ is \emph{Diophantine} if there are polynomials
$P_1,\dots,P_m\in\Z[k,x_1,\dots,x_u]$ with
\[
  k\in S \iff \exists x_1,\dots,x_u\in\N:\; P_1=\dots=P_m=0 .
\]
We call the pair $(u+1,\max_i\deg P_i)$ the \emph{cost} of the representation:
$u$ unknowns plus the parameter, and the largest degree occurring.
\end{definition}

\begin{definition}[Generator]
A polynomial $Q\in\Z[y_1,\dots,y_\nu]$ \emph{generates} $S\subseteq\N$ if
\[
  S=\{\,Q(y_1,\dots,y_\nu) : y_i\in\N \,\}\cap\Z_{>0},
\]
that is, the set of \emph{positive} values of $Q$ over nonnegative integer
arguments is exactly $S$.
\end{definition}

\begin{proposition}[Representation to generator]
\label{prop:gen}
If $S$ is represented as above, then
\[
  Q \;=\; (k+c)\cdot\Bigl(1-\sum_{i=1}^m P_i^{\,2}\Bigr)
\]
generates $S$ for the appropriate shift $c$, and
\[
  \deg Q \;=\; 1+2\max_i \deg P_i .
\]
\end{proposition}

\begin{proof}
The value of $Q$ is positive precisely when $\sum_i P_i^2=0$, i.e.\ when all
$P_i$ vanish, in which case $Q=k+c$. For the degree, the leading form of
$\sum_i P_i^2$ is a sum of squares of nonzero real polynomials and hence cannot
vanish; so $\deg\bigl(\sum_i P_i^2\bigr)=2\max_i\deg P_i$ exactly, and
multiplying by the linear factor adds one.
\end{proof}

\begin{remark}
\label{rem:odd}
Every generator \emph{obtained by Proposition~\ref{prop:gen}} has odd degree.
This is not a property of generators in general---$Q(y)=y^2$ generates
$\{1,4,9,\dots\}$ in the sense of Definition~2.2 and has degree $2$---but every
figure quoted in this literature comes from that construction, so parity is a
useful heuristic for deciding whether a quoted number is a representation degree
or a generator degree. We use it that way in \S\ref{ssec:parity}, and only that
way.
\end{remark}

\subsection{The two transformations}

\begin{definition}[Flattening]
Given an equation containing a subexpression $E$ with $\deg E\ge 2$, replace
every occurrence of $E$ by a fresh unknown $\mu$ and adjoin the equation
$\mu=E$. Iterating until every equation has degree at most $t$ yields a system
of degree $t$ with one extra unknown per name introduced.
\end{definition}

Flattening is always possible down to $t=2$, and $t=2$ is the floor: by
Proposition~\ref{prop:gen} a generator of degree~$3$ would require a
representation by linear equations, whose solution sets are semilinear and
cannot be the primes. Hence \emph{degree $5$ is the minimum attainable degree
for a prime generator by this route}, and the entire game in that corner is the
number of variables.

\begin{definition}[Elimination]
If some equation has the form $u=E$ with $u\notin\mathrm{vars}(E)$, then $u$ may
be substituted throughout and both the unknown and its equation dropped.
\end{definition}

Over $\N$ elimination is not unconditionally sound, and this is the point where
care is needed.

\begin{proposition}[Soundness of elimination over $\N$]
\label{prop:elim}
Let $\Sigma$ be a system over $\N$ containing $u=E$ with
$u\notin\mathrm{vars}(E)$, and let $\Sigma'$ be the result of substituting $E$
for $u$ and deleting that equation. If every solution of $\Sigma'$ satisfies
$E\ge 0$, then the solutions of $\Sigma$ project bijectively onto those of
$\Sigma'$; in particular the two systems are equisatisfiable.
\end{proposition}

\begin{proof}
A solution of $\Sigma$ restricts to one of $\Sigma'$. Conversely, given a
solution of $\Sigma'$, setting $u:=E$ satisfies $\Sigma$ provided $E\in\N$,
which is the hypothesis; the two maps are mutually inverse.
\end{proof}

\begin{remark}
The converse fails, and it is worth saying so because the criterion below is
applied mechanically. Equisatisfiability is a statement about the
\emph{existence} of solutions: $\Sigma'$ may have spurious solutions with $E<0$
and still be satisfiable exactly when $\Sigma$ is, because other solutions
survive. What the hypothesis buys is the stronger bijection, which is what one
needs in order to keep composing transformations.
\end{remark}

\begin{remark}[The criterion we use, and its limits]
\label{rem:crit}
A \emph{sufficient} condition for the hypothesis of
Proposition~\ref{prop:elim} is that every coefficient of $E$ be nonnegative, in
which case $E\ge 0$ holds on all of $\N^{u}$ for purely syntactic reasons. This
is the criterion our tooling applies by default, and it is decidable and cheap.
It is emphatically \emph{not necessary}: an expression such as $(A^2-1)C^2+1$
is nonnegative whenever $A\ge1$, which is a fact about the solution set rather
than about the shape of the polynomial. Whenever we eliminate on such grounds we
say so explicitly and prove the required bound; Sections~\ref{sec:var21}
and~\ref{sec:sec3} are largely about these proofs.
\end{remark}

\subsection{Domination}

Pairs $(\nu_1,\delta_1)$ and $(\nu_2,\delta_2)$ are compared componentwise: the
first \emph{dominates} the second if $\nu_1\le\nu_2$ and $\delta_1\le\delta_2$
with at least one inequality strict. A pair is a contribution when no known pair
dominates it. We use this criterion literally, including against our own
figures, and it is why some of the constructions below are reported as
superseded.
